\documentclass[10pt,letterpaper]{article}

% ---------- Page setup ----------
\usepackage[margin=0.7in]{geometry}
\usepackage[T1]{fontenc}
\usepackage[utf8]{inputenc}
\usepackage{charter}                 % Charter body
\usepackage[scaled=0.92]{helvet}     % Helvetica sans (headings)
\usepackage{microtype}
\usepackage{array}
\usepackage{booktabs}
\usepackage{tabularx}
\usepackage{longtable}
\usepackage{enumitem}
\usepackage{xcolor}
\usepackage{tikz}
\usetikzlibrary{arrows.meta,positioning,fit,backgrounds}
\usepackage[most]{tcolorbox}
\usepackage{titlesec}
\usepackage{fancyhdr}
\usepackage{lastpage}
\usepackage[hidelinks]{hyperref}
\usepackage{ragged2e}
\usepackage{multicol}

% ---------- Colors (house palette) ----------
\definecolor{Navy}{HTML}{17324D}
\definecolor{Blue}{HTML}{2563EB}
\definecolor{Teal}{HTML}{0F766E}
\definecolor{Green}{HTML}{15803D}
\definecolor{Amber}{HTML}{B45309}
\definecolor{Brick}{HTML}{B91C1C}
\definecolor{Purple}{HTML}{6D28D9}
\definecolor{GrayText}{HTML}{475569}
\definecolor{LightBlue}{HTML}{EFF6FF}
\definecolor{LightTeal}{HTML}{ECFDF5}
\definecolor{LightAmber}{HTML}{FFFBEB}
\definecolor{LightBrick}{HTML}{FEF2F2}
\definecolor{LightPurple}{HTML}{F5F3FF}
\definecolor{LightGreen}{HTML}{F0FDF4}
\definecolor{LightGray}{HTML}{F8FAFC}
\definecolor{RuleGray}{HTML}{CBD5E1}

% ---------- Global formatting ----------
\setlength{\parindent}{0pt}
\setlength{\parskip}{4pt}
\renewcommand{\arraystretch}{1.22}
\setlength{\tabcolsep}{4pt}
\setlist[itemize]{leftmargin=*, nosep, topsep=2pt}
\setlist[enumerate]{leftmargin=*, nosep, topsep=2pt}
\titleformat{\section}{\large\sffamily\bfseries\color{Navy}}{}{0em}{}[\titlerule]
\titleformat{\subsection}{\normalsize\sffamily\bfseries\color{Navy}}{}{0em}{}
\titleformat{\subsubsection}{\normalsize\sffamily\bfseries\color{Teal}}{}{0em}{}
\titlespacing*{\section}{0pt}{12pt}{5pt}
\titlespacing*{\subsection}{0pt}{9pt}{3pt}
\emergencystretch=2.5em
\sloppy

% ---------- Header/footer ----------
\setlength{\headheight}{14pt}
\pagestyle{fancy}
\fancyhf{}
\lhead{\sffamily\footnotesize\textcolor{GrayText}{Self-Hosting Cloud Architecture Reference}}
\rhead{\sffamily\footnotesize\textcolor{GrayText}{Enterprise AppSec + Operations Stack}}
\cfoot{\sffamily\footnotesize\textcolor{GrayText}{Page \thepage\ of \pageref{LastPage}}}
\renewcommand{\headrulewidth}{0.4pt}
\renewcommand{\footrulewidth}{0pt}

% ---------- Column types ----------
\newcolumntype{Y}{>{\RaggedRight\arraybackslash}X}
\newcolumntype{L}[1]{>{\RaggedRight\arraybackslash}p{#1}}
\newcolumntype{C}[1]{>{\Centering\arraybackslash}p{#1}}

% ---------- Boxes ----------
\newtcolorbox{summarybox}{colback=LightBlue, colframe=Blue, boxrule=0.7pt,
  arc=2mm, left=8pt, right=8pt, top=7pt, bottom=7pt}
\newtcolorbox{decisionbox}{colback=LightTeal, colframe=Teal, boxrule=0.7pt,
  arc=2mm, left=8pt, right=8pt, top=7pt, bottom=7pt}
\newtcolorbox{principlebox}{colback=LightTeal, colframe=Teal, boxrule=0.7pt,
  arc=2mm, left=8pt, right=8pt, top=7pt, bottom=7pt}
\newtcolorbox{warningbox}{colback=LightAmber, colframe=Amber, boxrule=0.7pt,
  arc=2mm, left=8pt, right=8pt, top=7pt, bottom=7pt}
\newtcolorbox{riskbox}{colback=LightBrick, colframe=Brick, boxrule=0.7pt,
  arc=2mm, left=8pt, right=8pt, top=7pt, bottom=7pt}

% ---------- Helpers ----------
\newcommand{\arch}[1]{\textcolor{Navy}{\textbf{#1}}}
\newcommand{\tool}[1]{\textbf{#1}}

\begin{document}

% ---------- Title ----------
{\sffamily\bfseries\color{Navy}\Huge Self-Hosting Cloud Architecture Reference}\par
\vspace{3pt}
{\sffamily\large\color{GrayText} Per-application deployment topology for a Kubernetes-centered, self-hosted Enterprise AppSec and Operations estate}\par
\vspace{4pt}
{\sffamily\small\color{GrayText} Companion to the Enterprise AppSec and Operations Stack and the Minimum Comprehensive Deployment Orchestration Stack \textbullet\ May 2026}\par
\vspace{10pt}

\begin{summarybox}
\textbf{Scope.} The deployment stack already fixes \tool{Kubernetes} as the runtime orchestrator, so the open question is not \emph{whether} to use Kubernetes but \emph{what topology each self-hosted component should adopt on or alongside that cluster}. This reference maps every self-hostable application in the stack to one of eight architecture archetypes, states the externalized state and high-availability posture for each, and identifies the cross-cutting cloud decisions --- externalized relational state, object-backed blob storage, and tier-based segmentation --- that do most of the work. One component, \tool{GHAS}, does not fit the containerized model and is handled as a dedicated appliance.
\end{summarybox}

\section{Purpose and Decision Standard}

This document answers a single question: for each application in the stack that must be self-hosted, what cloud architecture pattern is most appropriate? It does not re-litigate tool selection --- that is settled in the companion handouts --- and it does not enumerate every Kubernetes primitive. It assigns each component a deployment topology and the minimum set of design constraints required to make that topology safe to operate.

\begin{decisionbox}
\textbf{Decision rule.} The architecture for a component is determined by three properties, evaluated in order: (1) \emph{is it an application workload or a platform control plane?} (2) \emph{does it hold state, and can that state be externalized?} (3) \emph{does it scale out horizontally, or must it run as a low-replica or quorum-based singleton?} The answers place every component into exactly one of the archetypes below.
\end{decisionbox}

\section{The Deciding Axis: Statefulness and Scale-Out}

Most of these tools are described loosely as ``apps you can self-host,'' but they behave very differently at runtime. Three behaviors separate them:

\begin{itemize}
  \item \textbf{Twelve-factor services} (Authentik, Taiga, Harbor) externalize all state to a database, cache, and object store, so the running process is disposable and replicas scale horizontally.
  \item \textbf{Stateful monoliths} (Eramba, GLPI, iTop, SuiteCRM) assume local session and file state and run singleton cron, so they scale \emph{vertically} and depend on the database tier for durability rather than on replica count.
  \item \textbf{Control planes} (Kubernetes, Vault, Argo CD, admission and identity layers) are not application workloads at all; their failure cascades across every workload, so they are operated as managed services, quorum clusters, or highly-available in-cluster controllers.
\end{itemize}

The practical consequence is that ``put it on Kubernetes'' is correct for almost everything, but the \emph{shape} on Kubernetes --- horizontal autoscaling versus a pinned single replica, in-cluster versus externalized state, workload tenant versus management-plane tenant --- changes per component.

\section{Architecture Archetypes}

\begin{longtable}{L{0.55cm} L{3.6cm} L{6.2cm} L{4.3cm}}
\toprule
\textbf{} & \textbf{Archetype} & \textbf{Definition} & \textbf{Applies to} \\
\midrule
\endfirsthead
\toprule
\textbf{} & \textbf{Archetype} & \textbf{Definition} & \textbf{Applies to} \\
\midrule
\endhead
\bottomrule
\endfoot
\arch{A} & Managed control plane & Provider-operated Kubernetes control plane; the enterprise operates only node pools and add-ons. & Kubernetes itself. \\
\arch{B} & Stateless service + external data & Twelve-factor app tier (\texttt{Deployment} with optional autoscaling); database, cache, and blobs externalized off-cluster. & Authentik, Taiga, Moodle, Harbor, DongTai server. \\
\arch{C} & Single-instance + persistent volume & One replica (or active--passive) with a persistent volume for local file state, an external managed database, and a singleton cron. & Eramba, GLPI, iTop, SuiteCRM. \\
\arch{D} & HA stateful cluster & Quorum-based multi-node cluster with anti-affinity and durable storage; treated as tier-0 for recovery. & Vault; the long-term metrics and log tier. \\
\arch{E} & In-cluster operator / admission controller & A reconciling controller or admission webhook, run multi-replica for control-plane availability. & Argo CD, Argo Rollouts, External Secrets Operator, Kyverno / Gatekeeper, Trivy Operator. \\
\arch{F} & Ephemeral job / pipeline & Run-to-completion workload that scales to zero and leaves no persistent footprint. & Trivy (CI), OWASP ZAP, GitHub Actions runners, Cosign, Ansible. \\
\arch{G} & Agent / sidecar & Instrumentation attached to a target workload; never deployed standalone. & DongTai-IAST agent. \\
\arch{H} & Dedicated VM appliance & Vendor appliance with its own lifecycle, kept off the workload cluster. & GitHub Enterprise Server (GHAS). \\
\end{longtable}

\section{Business and Operations Applications}

\small
\begin{longtable}{L{1.15in} C{0.5in} L{4.9in}}
\toprule
\textbf{Application (role)} & \textbf{Arch.} & \textbf{Architecture and key design notes} \\
\midrule
\endfirsthead
\toprule
\textbf{Application (role)} & \textbf{Arch.} & \textbf{Architecture and key design notes} \\
\midrule
\endhead
\bottomrule
\endfoot
\tool{Authentik} \newline (IAM) & B & Server and worker \texttt{Deployment}s with external managed PostgreSQL and Redis. This is the identity control plane: give it the highest availability target in the estate and isolate its blast radius (treat as tier-0). \\
\tool{Taiga} \newline (Project mgmt) & B & Separate back, front, events, and async services with external PostgreSQL and a real RabbitMQ broker. The frontend is static (serve via ingress or CDN). Do not run the broker as a single fragile pod. \\
\tool{Moodle} \newline (LMS) & B & The one PHP application that genuinely scales out --- but only with \texttt{moodledata} on shared storage (EFS / Filestore or an object-storage plugin), sessions in Redis, and an external managed MySQL. Cron must be a single \texttt{CronJob}, never per-replica. \\
\tool{Eramba} \newline (GRC) & C & CakePHP monolith with singleton cron. Run a single replica with a persistent volume and an external managed MySQL / MariaDB. Scale vertically; rely on managed-database HA and point-in-time recovery rather than horizontal replicas. \\
\tool{GLPI} \newline (ITAM) & C & Low-replica web tier with a persistent volume for local file storage and an external MariaDB. Local uploads and the cron singleton make true scale-out risky. \\
\tool{iTop} \newline (ITSM / CMDB) & C & One to two replicas with an external MariaDB and a persistent volume. CMDB integrity matters most --- prioritize database point-in-time recovery over scale. \\
\tool{SuiteCRM} \newline (CRM) & C & Single or low-replica web tier with a persistent volume for uploads and an external MariaDB. Legacy PHP with local session and file assumptions; treat as a vertically scaled monolith with managed-database durability. \\
\tool{DongTai-IAST} \newline (IAST) & B + G & Split component. The \emph{server} (web, OpenAPI, engine) is a stateless tier with external MySQL and Redis. The \emph{agent} is archetype G: it instruments the target application pods as a sidecar or in-process agent and is never deployed standalone. \\
\tool{OWASP ZAP} \newline (DAST) & F & On-demand headless \texttt{Job} run against pre-production. Scale to zero; emit findings as release evidence, then discard the container. No persistent footprint. \\
\tool{GHAS} \newline (SAST) & H & Not a container workload. Self-hosted means GitHub Enterprise Server --- a dedicated VM / OVA appliance with its own backup and HA story, kept off the workload cluster. GitHub Enterprise Cloud is the managed (not self-hosted) alternative. See the dedicated section below. \\
\end{longtable}
\normalsize

\section{Platform and Deployment Components}

\small
\begin{longtable}{L{1.15in} C{0.5in} L{4.9in}}
\toprule
\textbf{Component} & \textbf{Arch.} & \textbf{Architecture and key design notes} \\
\midrule
\endfirsthead
\toprule
\textbf{Component} & \textbf{Arch.} & \textbf{Architecture and key design notes} \\
\midrule
\endhead
\bottomrule
\endfoot
\tool{Kubernetes} & A & Use a managed control plane (EKS / GKE / AKS) unless an air-gap or data-sovereignty constraint forbids it. The control plane is the one piece worth offloading to a provider; self-manage (kubeadm / Talos) only when required. \\
\tool{Harbor} \newline (registry) & B & Stateless components with external PostgreSQL and Redis, and \emph{object storage} (S3 / GCS) for image layers. Object-backed storage is what makes the registry cloud-appropriate; avoid block-volume image storage. \\
\tool{Argo CD} / \tool{Argo Rollouts} & E & In-cluster, HA mode (multiple repo-server and application-controller replicas, redis-ha). Run in a dedicated namespace, or on a separate management cluster, since it is itself a delivery control plane. \\
\tool{Vault} + \tool{External Secrets Operator} & D + E & Vault is archetype D: integrated Raft storage across three or five nodes, pod anti-affinity, and cloud-KMS auto-unseal --- the most operationally sensitive component, treated as tier-0 for availability and recovery. The External Secrets Operator is archetype E: an in-cluster operator that syncs external secrets into Kubernetes \texttt{Secret}s. \\
\tool{Kyverno} / \tool{OPA Gatekeeper} & E & Admission controllers run with multiple webhook replicas. A single-replica admission webhook is a cluster-wide availability risk, since its failure can block all admissions. \\
\tool{Prometheus} / \tool{Alertmanager} / \tool{Grafana} / \tool{Loki} & E + D & Collection runs in-cluster (archetype E); durable state is externalized (archetype D). Prometheus \texttt{remote\_write} to Thanos / Mimir or a managed store; Loki on object storage; Grafana on an external database; Alertmanager as a three-replica gossip mesh. Local volumes are acceptable only for short retention. \\
\tool{GitHub Actions} runners & F & Actions Runner Controller (ARC) with ephemeral pods that scale to zero. Isolate runner nodes from production workloads. \\
\tool{Trivy} & F + E & Dual mode: an ephemeral CI \texttt{Job} that gates registry promotion, and the Trivy Operator (archetype E) for continuous in-cluster CVE posture. \\
\tool{Cosign} / \tool{Sigstore} & F & Signing happens in CI (ephemeral); verification happens at admission via Kyverno / Gatekeeper policy. A standing service is required only if a Rekor / Fulcio transparency stack is self-hosted. \\
\tool{Ansible} & F & Procedural automation from a control node or CI runner --- not a persistent service. Introduce AWX on Kubernetes only if a self-service control layer is wanted. \\
\end{longtable}
\normalsize

\section{Reference Topology}

The estate resolves to four planes: a managed control plane, a tier-0 platform plane, application workloads, and ephemeral jobs --- all drawing on an off-cluster managed data plane, with the GHAS appliance deliberately detached.

\begin{center}
\begin{tikzpicture}[
  font=\sffamily\scriptsize,
  band/.style={rectangle, rounded corners=4pt, draw=#1, thick, fill=#1!6, minimum width=15.4cm},
  title/.style={anchor=west, font=\sffamily\bfseries\scriptsize, text=#1},
  box/.style={rectangle, rounded corners=3pt, draw=RuleGray, thick, fill=white,
    align=center, minimum height=0.58cm, text width=2.5cm, inner sep=2pt, font=\sffamily\scriptsize},
  arrow/.style={-{Latex[length=1.8mm]}, semithick, draw=GrayText}
]

% ---- Band A: Managed control plane ----
\node[band=Blue, minimum height=1.3cm] at (0,0) {};
\node[title=Blue] at (-7.5,0.45) {A \textbullet\ Managed control plane};
\node[box, draw=Blue, fill=LightBlue, text width=9cm] at (0,-0.12)
  {\textbf{Managed Kubernetes Control Plane} --- EKS / GKE / AKS (managed node pools and add-ons only)};

% ---- Band B: Tier-0 platform plane ----
\node[band=Purple, minimum height=1.55cm] at (0,-2.15) {};
\node[title=Purple] at (-7.5,-1.55) {B \textbullet\ Tier-0 platform plane (separate management cluster recommended)};
\node[box, draw=Purple] at (-5.7,-2.4) {\textbf{Authentik}\\Identity / SSO};
\node[box, draw=Purple] at (-2.85,-2.4) {\textbf{Vault} (HA Raft)\\+ ESO};
\node[box, draw=Purple] at (0,-2.4) {\textbf{Argo CD}\\/ Rollouts};
\node[box, draw=Purple] at (2.85,-2.4) {\textbf{Kyverno /}\\\textbf{Gatekeeper}};
\node[box, draw=Purple] at (5.7,-2.4) {\textbf{Prometheus}\\long-term tier};

% ---- Band C: Application workloads ----
\node[band=Teal, minimum height=2.7cm] at (0,-4.95) {};
\node[title=Teal] at (-7.5,-3.85) {C \textbullet\ Application workloads};
\node[box, draw=Teal] at (-5.7,-4.55) {\textbf{Taiga}};
\node[box, draw=Teal] at (-2.85,-4.55) {\textbf{Moodle}};
\node[box, draw=Teal] at (0,-4.55) {\textbf{Harbor}};
\node[box, draw=Teal] at (2.85,-4.55) {\textbf{DongTai} server};
\node[box, draw=Teal] at (5.7,-4.55) {\textbf{Grafana /}\\\textbf{Loki} read};
\node[box, draw=Teal] at (-5.7,-5.45) {\textbf{Eramba}};
\node[box, draw=Teal] at (-2.85,-5.45) {\textbf{GLPI}};
\node[box, draw=Teal] at (0,-5.45) {\textbf{iTop}};
\node[box, draw=Teal] at (2.85,-5.45) {\textbf{SuiteCRM}};
\node[box, draw=Teal] at (5.7,-5.45) {\textbf{DongTai} agent\\(sidecar)};

% ---- Band F: Ephemeral jobs ----
\node[band=GrayText, minimum height=1.7cm] at (0,-7.55) {};
\node[title=GrayText] at (-7.5,-6.9) {F \textbullet\ Ephemeral jobs / pipeline (scale to zero)};
\node[box, draw=GrayText] at (-5.7,-7.8) {\textbf{Trivy} (CI)};
\node[box, draw=GrayText] at (-2.85,-7.8) {\textbf{OWASP ZAP}};
\node[box, draw=GrayText] at (0,-7.8) {\textbf{Actions} runners};
\node[box, draw=GrayText] at (2.85,-7.8) {\textbf{Cosign}};
\node[box, draw=GrayText] at (5.7,-7.8) {\textbf{Ansible}};

% ---- Band D: External managed data plane ----
\node[band=Green, minimum height=1.7cm] at (0,-9.65) {};
\node[title=Green] at (-7.5,-9.0) {D \textbullet\ External managed data plane (off-cluster)};
\node[box, draw=Green, fill=LightGreen] at (-5.7,-9.9) {Managed SQL\\(Postgres / MySQL)};
\node[box, draw=Green, fill=LightGreen] at (-2.85,-9.9) {Managed Redis\\/ broker};
\node[box, draw=Green, fill=LightGreen] at (0,-9.9) {Object storage\\(S3 / GCS)};
\node[box, draw=Green, fill=LightGreen] at (2.85,-9.9) {Cloud KMS\\(auto-unseal)};
\node[box, draw=Green, fill=LightGreen] at (5.7,-9.9) {Long-term metrics\\(Thanos / Mimir)};

% ---- Detached GHAS appliance ----
\node[box, draw=Brick, fill=LightBrick, text width=8.4cm] (ghas) at (0,-11.2)
  {\textbf{H \textbullet\ GitHub Enterprise Server (GHAS)} --- dedicated off-cluster appliance};

% ---- Arrows ----
% State externalization: application workloads -> external data plane (left channel)
\draw[arrow] (-7.7,-4.95) -- (-8.15,-4.95) -- (-8.15,-9.65) -- (-7.7,-9.65);
\node[anchor=south, rotate=90, font=\sffamily\scriptsize, text=GrayText] at (-8.4,-7.3) {externalized state};
% GHAS feeds the pipeline but stays detached (right channel, dashed)
\draw[arrow, draw=Brick, dashed] (4.2,-11.2) -- (8.15,-11.2) -- (8.15,-7.55) -- (7.7,-7.55);
\node[anchor=north, rotate=90, font=\sffamily\scriptsize, text=Brick] at (8.4,-9.4) {detached feed};
\end{tikzpicture}
\end{center}

\vspace{2pt}
{\footnotesize\color{GrayText} Solid arrows: application and platform planes draw their relational, cache, and blob state from the off-cluster managed data plane. Dashed arrow: GHAS feeds source-security results into the pipeline but is operated as a detached appliance, not a cluster tenant.}

\section{Cross-Cutting Cloud Decisions}

Three decisions account for most of the resilience and scalability outcomes across every row above. They apply regardless of which archetype a component sits in.

\begin{principlebox}
\textbf{1. Externalize all relational state.} Eramba, GLPI, iTop, Moodle, SuiteCRM, Authentik, Taiga, Harbor, and the DongTai server all want their database \emph{out} of the cluster on managed MySQL or PostgreSQL (RDS, Cloud SQL, Aurora). This is what turns a stateful monolith into a restartable pod, and it moves high-availability, backup, and point-in-time recovery onto a service built to do them well. Avoid running production databases as in-cluster \texttt{StatefulSet}s unless no managed option exists.
\end{principlebox}

\begin{principlebox}
\textbf{2. Push blob and file state to object storage.} Harbor image layers, Loki log chunks, Moodle's \texttt{moodledata}, and application file uploads belong on S3 / GCS --- or on a shared filesystem (EFS / Filestore) where object storage is unsupported. This is the difference between a replica that scales freely and a replica pinned to one node's disk.
\end{principlebox}

\begin{principlebox}
\textbf{3. Segment by tier, not just by namespace.} The identity plane (Authentik), the secrets plane (Vault), and the policy plane (Kyverno / Gatekeeper) are tier-0: their failure cascades across every workload. Running these on a separate management or platform cluster --- distinct from application workloads --- is worth the operational overhead given this stack's compliance posture.
\end{principlebox}

\section{Availability and Recovery Tiers}

The archetypes also map to a recovery-priority model, which is the lens that matters most during incident planning and disaster-recovery design.

\begin{tabularx}{\textwidth}{L{1.5in} L{2.0in} Y}
\toprule
\textbf{Tier} & \textbf{Components} & \textbf{Why this tier} \\
\midrule
Tier-0 (platform / control plane) & Kubernetes control plane, Authentik, Vault, Argo CD, Kyverno / Gatekeeper, External Secrets Operator & Failure cascades across all workloads; warrants HA topology, anti-affinity, and a separate management cluster. \\
Tier-1 (systems of record) & Eramba, iTop, GLPI, Harbor & Data integrity and point-in-time recovery are the dominant concern; managed-database durability and tested restores are mandatory. \\
Tier-2 (delivery and enablement) & Taiga, Moodle, SuiteCRM, DongTai server, observability read paths & Important but tolerant of short outages; standard backups and single or low-replica topologies are acceptable. \\
Ephemeral (no standing SLA) & Trivy, OWASP ZAP, Actions runners, Cosign, Ansible & Run-to-completion jobs with no persistent state; resilience is achieved by re-running, not by replication. \\
\bottomrule
\end{tabularx}

\section{Adoption Sequencing}

\begin{enumerate}
  \item Stand up the managed Kubernetes control plane and the off-cluster managed data plane (SQL, Redis, object storage, KMS) first; everything else assumes they exist.
  \item Establish the tier-0 platform plane --- Authentik, Vault with auto-unseal, Argo CD, and admission policy --- ideally on a dedicated management cluster.
  \item Onboard archetype-B services (Authentik already done; then Harbor, Taiga, Moodle, DongTai server) against the externalized data plane.
  \item Migrate the archetype-C monoliths (Eramba, GLPI, iTop, SuiteCRM) with persistent volumes for file state and managed databases for durability; convert cron to singleton \texttt{CronJob}s.
  \item Wire in archetype-F pipelines (Trivy, ZAP, Actions runners, Cosign, Ansible) and the DongTai agent against running workloads.
  \item Provision the GHAS appliance (GitHub Enterprise Server) on its own lifecycle and connect it to the source and CI control loops.
\end{enumerate}

\section{Special Case: GHAS / GitHub Advanced Security}

\begin{warningbox}
\textbf{GHAS is the only component that does not fit the model.} Everything else in the stack is a containerized workload or a run-to-completion job. GHAS is a commercial capability of GitHub, not an open-source platform you deploy as a pod. ``Self-hosted'' for GHAS means \emph{GitHub Enterprise Server} --- a dedicated virtual appliance (OVA) with its own backup-utilities, replica, and upgrade lifecycle. Plan it as a separate system with its own availability story, not as a tenant of the workload cluster. If a self-hosted requirement is not strictly mandatory, GitHub Enterprise Cloud delivers the same GHAS capabilities as a managed service and removes the appliance-operations burden entirely.
\end{warningbox}

\section{Recommended Decision}

\begin{riskbox}
\textbf{Final decision.} Run a \tool{managed Kubernetes control plane} with an \tool{off-cluster managed data plane} (SQL, Redis, object storage, KMS). Deploy \tool{Authentik}, \tool{Taiga}, \tool{Moodle}, \tool{Harbor}, and the \tool{DongTai server} as stateless tiers with externalized state (archetype B). Deploy \tool{Eramba}, \tool{GLPI}, \tool{iTop}, and \tool{SuiteCRM} as single-instance monoliths with persistent volumes and managed databases (archetype C). Operate \tool{Vault} as an HA Raft cluster and the long-term metrics tier as durable stores (archetype D); run \tool{Argo CD}, \tool{Argo Rollouts}, \tool{ESO}, \tool{Kyverno / Gatekeeper}, and the \tool{Trivy Operator} as HA in-cluster controllers (archetype E); run \tool{Trivy} CI scans, \tool{OWASP ZAP}, \tool{GitHub Actions} runners, \tool{Cosign}, and \tool{Ansible} as ephemeral jobs (archetype F); attach the \tool{DongTai agent} as a sidecar (archetype G). Operate \tool{GHAS} as a dedicated \tool{GitHub Enterprise Server} appliance (archetype H), detached from the workload cluster. Place all tier-0 planes on a separate management cluster.
\end{riskbox}

\vfill
\begin{center}
{\sffamily\footnotesize\color{GrayText} End of reference}
\end{center}

\end{document}
